% Appendix: core code + AI tool usage report
\lstset{
  language=Python,
  basicstyle=\ttfamily\footnotesize,
  keywordstyle=\bfseries,
  commentstyle=\itshape\color{gray},
  stringstyle=\color{blue!60!black},
  numbers=left,
  numberstyle=\tiny\color{gray},
  breaklines=true,
  showstringspaces=false,
  frame=single,
  framerule=0.3pt,
  columns=fullflexible,
  keepspaces=true,
  xleftmargin=1.2em,
}

\section*{Declaration of AI Tool Usage}
\addcontentsline{toc}{section}{Declaration of AI Tool Usage}

In accordance with the competition's rules on generative-AI use, we disclose that an
AI pair-programming assistant (built into the Qoder development environment,
\url{https://qoder.com}) was used in three bounded roles: (i) implementing, debugging
and regression-testing the Python code realizing our models (geometry predicates,
Dijkstra routing, Fisher information, simulator, scenario batching); (ii) polishing
the English wording of sections the team had already outlined, argued and numerically
filled; (iii) assisting with \LaTeX{} typesetting. \textbf{All modeling ideas,
mathematical derivations, convention decisions, optimality arguments, and every
numerical conclusion were produced, checked and approved by the team}; AI served as
an implementation and language tool, not a source of results. Three in-text
annotations mark substantive assistance (\S\ref{sec:analysis}*, \S\ref{sec:q4}*,
Appendix~A code headers carry AI-assistance comments), and the detailed usage record
is Appendix~B. We take full responsibility for the entire submission.

\section*{Appendix}
\addcontentsline{toc}{section}{Appendix}

\subsection*{A. Core Code}
\addcontentsline{toc}{subsection}{A. Core Code}

All code is Python~3. The listings below are the computational cores of each question,
condensed from the full implementation (comments translated into English); the complete
repository structure is summarized at the end of this appendix. Header comments mark
where AI assistance was used in implementation.

\begin{lstlisting}[caption={Shared geometry engine: closed-set three-constraint
predicate (AI-assisted implementation, team-specified conventions).},label={lst:geom}]
VEHICLE_RADIUS = 1.5; SPEED = 3.0; CAMERA_HEIGHT = 2.0
ELEV_MIN, ELEV_MAX = -10.0, 80.0        # deg, closed interval
EPS = 1e-9                              # closed-set tolerance

def dist_point_to_rect(px, py, x0, x1, y0, y1):
    dx = max(x0 - px, 0.0, px - x1)
    dy = max(y0 - py, 0.0, py - y1)
    return math.hypot(dx, dy)

def clearance_ok(p, margin=VEHICLE_RADIUS):
    # disc center feasible: 1.5 m from field edge and all buildings
    # (exactly margin -> admissible, closed set)
    x, y = p
    if not (margin - EPS <= x <= 100 - margin + EPS
            and margin - EPS <= y <= 100 - margin + EPS):
        return False
    for b in BUILDINGS:
        if dist_point_to_rect(x, y, b.x0, b.x1, b.y0, b.y1) < margin - EPS:
            return False
    return True

def segment_hits_box(p0, p1, box, eps=EPS):
    # 3D slab test: sightline vs closed prism (tangency = blocked)
    xmin, xmax, ymin, ymax, zmin, zmax = box
    d = (p1[0]-p0[0], p1[1]-p0[1], p1[2]-p0[2])
    tmin, tmax = 0.0, 1.0
    for a, (lo, hi) in enumerate(((xmin,xmax),(ymin,ymax),(zmin,zmax))):
        if abs(d[a]) < 1e-15:
            if p0[a] < lo - eps or p0[a] > hi + eps: return False
        else:
            t1, t2 = (lo-p0[a])/d[a], (hi-p0[a])/d[a]
            if t1 > t2: t1, t2 = t2, t1
            tmin, tmax = max(tmin, t1), min(tmax, t2)
            if tmin > tmax + eps: return False
    return True

def los_ok(p, target):                  # camera at disc center, height 2 m
    p0 = (p[0], p[1], CAMERA_HEIGHT)
    p1 = (target.x, target.y, target.H)
    return not any(segment_hits_box(p0, p1,
                   (b.x0, b.x1, b.y0, b.y1, 0.0, b.h)) for b in BUILDINGS)

def elevation_ok(p, target):
    r = math.hypot(target.x - p[0], target.y - p[1])
    e = math.degrees(math.atan2(target.H - CAMERA_HEIGHT, r))
    return ELEV_MIN - EPS <= e <= ELEV_MAX + EPS

def is_station_feasible(p, target):
    return clearance_ok(p) and elevation_ok(p, target) and los_ok(p, target)
\end{lstlisting}

\begin{lstlisting}[caption={Exact disc-clearance for a travel segment and the
visibility-graph router (AI-assisted implementation).},label={lst:path}]
def segment_rect_distance(p, q, x0, x1, y0, y1):
    # exact min distance segment->rectangle (0 if entering);
    # inflated rectangles would wrongly reject legal corner-grazing paths
    if segment_crosses_rect_interior(p, q, (x0, x1, y0, y1)):
        return 0.0
    return min(dist_point_to_rect(*p, x0, x1, y0, y1),
               dist_point_to_rect(*q, x0, x1, y0, y1),
               *(dist_point_to_segment(cx, cy, p, q)
                 for cx in (x0, x1) for cy in (y0, y1)))

def shortest_path(start, goal, margin=VEHICLE_RADIUS):
    # corner visibility graph over inflated buildings + Dijkstra
    if not (clearance_ok(start, margin) and clearance_ok(goal, margin)):
        raise ValueError("endpoint violates vehicle geometry")
    rects = [inflate(b) for b in BUILDINGS]      # b grown by margin
    nodes = [start, goal] + [corners(r) for r in rects]
    adj = pairwise_edges(nodes)   # edge kept iff every
    ...           # segment_rect_distance(node_i, node_j, b) >= margin
    return dijkstra(adj, 0, 1)    # (length, path); None if unreachable
\end{lstlisting}

\begin{lstlisting}[caption={Q2(1): Fisher information, RMSE closed form and
attainable lower bound (derivation by the team, implementation AI-assisted).},
label={lst:fisher}]
def fisher_rmse(s1, s2, Tp, sigma_b):
    # I = sigma_b^-2 * sum(n_i n_i^T / d_i^2);  RMSE = sqrt(tr(C))
    n = [(Tp - s) / norm(Tp - s) for s in (s1, s2)]
    d = [norm(Tp - s) for s in (s1, s2)]
    I = sum(np.outer(ni, ni) / di**2 for ni, di in zip(n, d)) / sigma_b**2
    C = np.linalg.inv(I)
    return math.sqrt(np.trace(C))          # = sigma_b*sqrt(d1^2+d2^2)/sin(gamma)

def lower_bound(sigma_b, r80):
    return math.sqrt(2.0) * sigma_b * r80  # attained iff gamma=90, d1=d2=r80
\end{lstlisting}

\begin{lstlisting}[caption={Q2(2)/Q4: event-timeline simulator core (joint start =
max readiness; shared by both questions, replayed from captured states in Q4;
AI-assisted implementation).},label={lst:sim}]
def simulate(tasks, t0=0.0):
    # tasks: per vehicle, ordered (alert, station) items
    events = {v: [] for v in tasks}
    ready = {}                              # alert -> [ready_1, ready_2]
    for v, seq in tasks.items():
        t = t0
        pos, pan, tilt = INIT[v], 0.0, 0.0
        for a, st in seq:
            L, path = shortest_path(pos, st)        # feasible upper bound
            events[v].append(("arrive", a, t, t + L / SPEED)); t += L / SPEED
            s = slew_seconds(pan, tilt, pan_of(st, a), tilt_of(st, a))
            events[v].append(("slew", a, t, t + s)); t += s
            ready.setdefault(a, []).append(t)
            obs_start = max(ready[a])               # joint start
            events[v].append(("obs", a, obs_start, obs_start + 5.0))
            t = obs_start + 5.0
            r = reset_seconds(pan_of(st, a), tilt_of(st, a))
            events[v].append(("reset", a, t, t + r)); t += r
            pos, pan, tilt = st, pan_of(st, a), tilt_of(st, a)
    return events
\end{lstlisting}

\begin{lstlisting}[caption={Q3: pointing-error propagation chain and the inverse
$s_{\max}$ (AI-assisted implementation).},label={lst:q3}]
def sigma_tot(s, a=0.3, b=0.02, r=40.0):
    sigma_loc = a + b * s
    sigma1  = (180.0/math.pi) * sigma_loc / r     # deg, small-angle
    return math.sqrt(sigma1**2 + 0.8**2 + 0.6**2)

def p_reliable(s):                                # P(|dpsi| <= 5 deg)
    return 2.0 * Phi(5.0 / sigma_tot(s)) - 1.0

def s_max():                                      # invert the 95% chain
    z = 1.959964                                  # Phi^-1(0.975)
    sigma_tot_cap = 5.0 / z                       # 2.551067 deg
    sigma1_cap = math.sqrt(sigma_tot_cap**2 - 0.8**2 - 0.6**2)
    sigma_loc_cap = sigma1_cap * 40.0 / (180.0/math.pi)   # 1.638446 m
    return (sigma_loc_cap - 0.3) / 0.02           # 66.9223 m
\end{lstlisting}

\begin{lstlisting}[caption={Q4: rule-based strategy P2 (earliest-completion greedy;
AI-assisted implementation, rule designed by the team).},label={lst:p2}]
def strategy_P2(open_alerts, states, tau):
    ready = {}                       # alert -> start time once both assigned
    free  = {v: st for v, st in states.items()}   # pos, attitude, free_t
    remaining = set(open_alerts)
    while remaining:
        best = None
        for a in remaining:          # earliest achievable completion
            for v1, v2 in combinations(free, 2):
                if azimuth_gap_ok(v1, v2, a):
                    t = joint_ready(v1, v2, a, free, tau)
                    if best is None or t < best[0]:
                        best = (t, a, v1, v2)
        if best is None: return None          # no feasible assignment
        t, a, v1, v2 = best
        ready[a] = t; remaining.discard(a)
        for v in (v1, v2): free.pop(v)        # busy until task done
    return lexicographic_score(ready)         # (k, then -t_last)
\end{lstlisting}

\textbf{Repository structure.} \texttt{geometry.py} (engine, Listings~1--2),
\texttt{q1\_observable\_domain.py} (Q1 grid + analytic arc audit),
\texttt{q2\_1\_localization.py} (Q2(1) Fisher model, Listings~3),
\texttt{q2\_2\_scheduler.py} (Q2(2) 360-structure search, Listing~4),
\texttt{q3\_error\_propagation.py} (Q3 chain, Listing~5),
\texttt{q4\_rescheduler.py} (Q4 scenarios and strategies, Listing~6),
\texttt{simulator.py}, \texttt{route\_audit.py} (independent route legality audit),
\texttt{test\_modules.py} (51-check regression suite) and figure scripts; all random
components use fixed seed 2026 and every reported number and figure regenerates from
them.

\subsection*{B. AI Tool Usage Report}
\addcontentsline{toc}{subsection}{B. AI Tool Usage Report}

\noindent\textbf{Tool.} Qoder Desktop built-in AI pair-programming assistant
(Qoder, \url{https://qoder.com}), used during 2026-11-20/21 in the team's development
environment.

\medskip
\noindent\textbf{Scope of use.}
\begin{enumerate}
\item \textbf{Code implementation and debugging.} Translating the team's
model specifications into Python (geometry predicates, visibility-graph routing,
Fisher-information computation, 360-structure scheduler, error-propagation chain,
86-scenario rescheduler), fixing numerical defects found by tests (e.g.\ a
sharp-corner clearance over-approximation in the inflated-rectangle edge test), and
running the 51-check regression suite.
\item \textbf{Verification tooling.} Implementing the exact-geometry Monte Carlo
replays (seed 2026), the independent route-leg auditor, and the scenario batch
evaluation used in Q4.
\item \textbf{Language and typesetting.} Polishing English wording of sections whose
content, structure, derivations and numbers the team supplied; \LaTeX{} formatting
against the official template.
\end{enumerate}

\medskip
\noindent\textbf{What the AI did not do.} The AI did not originate the modeling
approach, the convention interpretations of the problem statement, the mathematical
derivations (observable-region arcs, Fisher lower bound and its attainment condition,
min--max structure enumeration, error-propagation chain), the optimality and
validation arguments, or the conclusions. All such content was specified, verified
numerically, and approved by the team; every reported number traces to a team-run
script listed in Appendix~A.

\medskip
\noindent\textbf{Usage records.} Sessions are logged in the team's shared repository
(commit history and analysis notes); three substantive in-text annotations appear at
\S2.1 (footnote), \S4.5 (footnote), and the Appendix~A listing headers, consistent
with this report and with the Declaration at the head of the Appendix.

\medskip
\noindent\textbf{Responsibility.} The team takes full responsibility for the validity
of the models, the correctness of the code, and all conclusions of this paper.
